% basics
\usepackage{iftex}
\ifXeTeX
  \usepackage{fontspec}
  \usepackage{xeCJK}
  \IfFontExistsTF{Noto Serif CJK TC}{%
    \setCJKmainfont{Noto Serif CJK TC}[BoldFont=Noto Serif CJK TC Bold]%
    \setCJKsansfont{Noto Sans CJK TC}[BoldFont=Noto Sans CJK TC Bold]%
    \setCJKmonofont{Noto Sans Mono CJK TC}%
  }{% Fallback to macOS system CJK fonts when Noto CJK is not installed
    \setCJKmainfont{Songti TC}%
    \setCJKsansfont{PingFang TC}%
    \setCJKmonofont{PingFang TC}%
  }
\else\ifLuaTeX
  \usepackage{luatexja-fontspec}
  \IfFontExistsTF{Noto Serif CJK TC}{%
    \setmainjfont{Noto Serif CJK TC}[BoldFont=Noto Serif CJK TC Bold]%
    \setsansjfont{Noto Sans CJK TC}[BoldFont=Noto Sans CJK TC Bold]%
    \setmonojfont{Noto Sans Mono CJK TC}%
  }{% Fallback to macOS system CJK fonts when Noto CJK is not installed
    \setmainjfont{Songti TC}%
    \setsansjfont{PingFang TC}%
    \setmonojfont{PingFang TC}%
  }
\else
  \PackageError{header}{This template requires XeLaTeX or LuaLaTeX}{}
\fi\fi
\IfFontExistsTF{Latin Modern Sans}{%
  \setsansfont{Latin Modern Sans}%
}{% Load Latin Modern Sans by file name (shipped with TeX Live, not always registered)
  \setsansfont{lmsans10-regular.otf}[%
    BoldFont       = lmsans10-bold.otf,%
    ItalicFont     = lmsans10-oblique.otf,%
    BoldItalicFont = lmsans10-boldoblique.otf,%
  ]%
}
\newfontface\lmrfont{lmroman10-bold.otf}

\usepackage[utf8]{inputenc}
\usepackage[T1]{fontenc}
\usepackage{textcomp}
\usepackage[hyphens]{url}
\usepackage[style=alphabetic,maxcitenames=1]{biblatex}
\usepackage[colorlinks=true,linkcolor=cyan,urlcolor=magenta,citecolor=violet]{hyperref}
\usepackage{graphicx}
\usepackage{float}
\usepackage{array}
\usepackage{booktabs}
% Local switch for a tabular whose `|' separators must run unbroken through the
% booktabs rules: kill the rule separation that breaks them, and give back the
% breathing room with extra row height.
\newcommand{\solidrules}{%
	\setlength{\aboverulesep}{0pt}%
	\setlength{\belowrulesep}{0pt}%
	\setlength{\extrarowheight}{0.75ex}%
}
\usepackage{emptypage}
\usepackage{subcaption}
\usepackage[usenames,dvipsnames,svgnames]{xcolor}
% Shaded rows and cells in the dichotomy tables (\rowcolor, \cellcolor).
\usepackage{colortbl}
% Row labels spanning several table rows (\multirow).
\usepackage{multirow}
% quiver style
\usepackage{tikz-cd}
% `calc` is necessary to draw curved arrows.
\usetikzlibrary{calc}
% `pathmorphing` is necessary to draw squiggly arrows.
\usetikzlibrary{decorations.pathmorphing}

% A TikZ style for curved arrows of a fixed height, due to AndréC.
\tikzset{curve/.style={settings={#1},to path={(\tikztostart)
					.. controls ($(\tikztostart)!\pv{pos}!(\tikztotarget)!\pv{height}!270:(\tikztotarget)$)
					and ($(\tikztostart)!1-\pv{pos}!(\tikztotarget)!\pv{height}!270:(\tikztotarget)$)
					.. (\tikztotarget)\tikztonodes}},
	settings/.code={\tikzset{quiver/.cd,#1}
			\def\pv##1{\pgfkeysvalueof{/tikz/quiver/##1}}},
	quiver/.cd,pos/.initial=0.35,height/.initial=0}

% TikZ arrowhead/tail styles.
\tikzset{tail reversed/.code={\pgfsetarrowsstart{tikzcd to}}}
\tikzset{2tail/.code={\pgfsetarrowsstart{Implies[reversed]}}}
\tikzset{2tail reversed/.code={\pgfsetarrowsstart{Implies}}}
% TikZ arrow styles.
\tikzset{no body/.style={/tikz/dash pattern=on 0 off 1mm}}

% useful macro for class
\newcommand{\at}[3]{\left.#1\right\vert_{#2}^{#3}}
\newcommand\quotient[2]{
	\mathchoice
	{% \displaystyle
		\text{\raise1ex\hbox{$#1$}\Big/\lower1ex\hbox{$#2$}}%
	}
	{% \textstyle
		#1\,/\,#2
	}
	{% \scriptstyle
		#1\,/\,#2
	}
	{% \scriptscriptstyle
		#1\,/\,#2
	}
}

\newcommand{\True}{\textsf{True}}
\newcommand{\T}{\textsf{T}}
\newcommand{\False}{\textsf{False}}
\newcommand{\F}{\textsf{F}}
\newcommand{\act}{\rotatebox[origin=c]{-180}{\(\,\circlearrowright\,\)}}

\usepackage{physics}
\usepackage{complexity}

\DeclareMathOperator{\id}{id}
\DeclareMathOperator{\Span}{span}
\DeclareMathOperator{\supp}{supp}
\DeclareMathOperator{\vol}{vol}
\DeclareMathOperator{\dist}{dist}
\DeclareMathOperator{\diam}{diam}
\DeclareMathOperator{\im}{Im}
\DeclareMathOperator{\sgn}{sgn}
\DeclareMathOperator{\sign}{sign}
\DeclareMathOperator{\Int}{Int}
\DeclareMathOperator{\diag}{diag}
\DeclareMathOperator{\dom}{dom}
\DeclareMathOperator*{\argmax}{arg\,max}
\DeclareMathOperator*{\argmin}{arg\,min}
% In-sample and out-of-sample error, from the feasibility chapter on.
\newcommand{\Ein}{E_{\mathrm{in}}}
\newcommand{\Eout}{E_{\mathrm{out}}}
% VC dimension, from the VC-dimension chapter on.
\newcommand{\dvc}{d_{\mathrm{VC}}}
% Pointwise error measure and its algorithmic surrogate (noise-and-error chapter).
\newcommand{\err}{\mathrm{err}}
\newcommand{\errhat}{\widehat{\mathrm{err}}}

\DeclareMathOperator{\Var}{Var}
\DeclareMathOperator{\Cov}{Cov}
\DeclareMathOperator{\Corr}{Corr}

\let\implies\Rightarrow
\let\impliedby\Leftarrow
\let\iff\Leftrightarrow

\usepackage{stmaryrd} % for \lightning
\newcommand\conta{\scalebox{1.1}{\(\lightning\)}}

\usepackage{bm}
\usepackage{bbm}

% figure support
\usepackage{import}
\usepackage{xifthen}
\ifdefined\pdfminorversion
  \pdfminorversion=7
\fi

\usepackage{pdfpages}
\usepackage{transparent}
\newcommand{\incfig}[2][\columnwidth]{%
	\def\svgwidth{#1}
	\import{Figures/}{#2.pdf_tex}
}

\makeatletter
\def\moverlay{\mathpalette\mov@rlay}
\def\mov@rlay#1#2{\leavevmode\vtop{%
   \baselineskip\z@skip \lineskiplimit-\maxdimen
   \ialign{\hfil$\m@th#1##$\hfil\cr#2\crcr}}}
\newcommand{\charfusion}[3][\mathord]{
    #1{\ifx#1\mathop\vphantom{#2}\fi
        \mathpalette\mov@rlay{#2\cr#3}
      }
    \ifx#1\mathop\expandafter\displaylimits\fi}
\makeatother

\newcommand{\cupdot}{\charfusion[\mathbin]{\cup}{\cdot}}
\newcommand{\bigcupdot}{\charfusion[\mathop]{\bigcup}{\cdot}}

\usepackage{amsmath, amsfonts, mathtools, amsthm, amssymb, mathrsfs}
\usepackage{geometry}
\usepackage{cancel}
\usepackage{systeme}

\usepackage[inline, shortlabels]{enumitem}
\usepackage{multicol}
\setlength\multicolsep{0pt}

\usepackage{caption}
\captionsetup{belowskip=0pt}
\geometry{a4paper,left=2.54cm,right=2.54cm,top=2.54cm,bottom=2.54cm}

% for the big braces
\usepackage{bigdelim}

\definecolor{correct}{HTML}{009900}
\newcommand\correct[2]{\ensuremath{\:}{\color{red}{#1}}\ensuremath{\to }{\color{correct}{#2}}\ensuremath{\:}}
\newcommand\green[1]{{\color{correct}{#1}}}
\newcommand\hide[1]{}

% si unitx
\usepackage{siunitx}
\sisetup{locale = FR}
\newcommand\mat[1]{\mathbf{#1}}

% tikz
\usepackage{tikz}
\usetikzlibrary{intersections, angles, quotes, positioning}
\usetikzlibrary{arrows.meta}
\usetikzlibrary{patterns}
\usepackage{pgfplots}
\pgfplotsset{compat=1.13}
\usepgfplotslibrary{groupplots}
\usepgfplotslibrary{fillbetween}

\tikzset{
	force/.style={thick, {Circle[length=2pt]}-stealth, shorten <=-1pt}
}

% Styles for the machine-learning flow diagrams:
%   mlbox   — a passive object (data, target function, hypothesis, ...)
%   mlop    — an active step that consumes and produces (learning, ML, A)
%   mlarrow — the flow between them
\tikzset{
	mlbox/.style={draw=black!65, rounded corners=3pt, inner sep=5pt,
			minimum height=9mm, minimum width=2.4cm, align=center, font=\small},
	mlop/.style={mlbox, fill=TealBlue!12, draw=TealBlue!60!black,
			font=\small\bfseries},
	mlarrow/.style={-{Stealth[length=2.2mm]}, thick, draw=black!65},
}

% Styles for classification pictures: a hollow circle marks a \(+1\) example and
% a cross a \(-1\) one, while the two shaded half-spaces show which side of a
% hypothesis is called positive.
\tikzset{
	dpos/.style={circle, draw=blue!80!black, semithick, inner sep=0pt,
			minimum size=4.5pt},
	dneg/.style={red!85!black, inner sep=0pt, font=\small},
	poshalf/.style={fill=blue!9},
	neghalf/.style={fill=red!9},
	sepline/.style={thick, draw=black!85},
	wvec/.style={-{Stealth[length=2.2mm]}, very thick, draw=Purple},
}

% Marks for multi-class / unlabelled scatter plots (coin recognition and any
% later multi-class illustration).  Usage: \node[dcls=ForestGreen] at (..) {};
\tikzset{
	dcls/.style={circle, draw=#1!85!black, fill=#1!35, semithick,
			inner sep=0pt, minimum size=3.4pt},
	dunl/.style={circle, draw=black!50, semithick,
			inner sep=0pt, minimum size=3.4pt},
	plotbox/.style={draw=black!55, semithick},
}

% The deck's own four-class scatter plots (multiclass sections of `Linear Models
% for Classification'): one marker shape per class --- mcsq (square), mcdia
% (diamond), mctri (triangle), mcstar (star) --- and one pale tint per class,
% rgsq / rgdia / rgtri / rgstar, for the decision regions of a combined
% multiclass classifier.  `tieregion' marks inputs on which the combination has
% no unique answer.  \mcmark{sq|dia|tri|star} gives the marker at text size,
% for captions such as `square or not?'.
\usetikzlibrary{shapes.geometric}
\definecolor{clsq}{HTML}{86AD86}
\definecolor{cldia}{HTML}{A3A374}
\definecolor{cltri}{HTML}{AD86AD}
\definecolor{clstar}{HTML}{86ADB8}
\colorlet{rgsq}{green!13}
\colorlet{rgdia}{yellow!24}
\colorlet{rgtri}{magenta!12}
\colorlet{rgstar}{cyan!13}
\tikzset{
	mcsq/.style={rectangle, fill=clsq, inner sep=0pt, minimum size=4.2pt},
	mcdia/.style={diamond, fill=cldia, inner sep=0pt, minimum width=4.6pt,
			minimum height=6.6pt},
	mctri/.style={regular polygon, regular polygon sides=3, fill=cltri,
			inner sep=0pt, minimum size=6.8pt},
	mcstar/.style={star, star points=5, star point ratio=2.3, fill=clstar,
			inner sep=0pt, minimum size=6.6pt},
	tieregion/.style={pattern=crosshatch dots, pattern color=black!40},
}
\newcommand{\mcmark}[1]{\tikz[baseline=-0.6ex]{\node[mc#1] at (0,0) {};}}

% Panels of the `how many kinds of lines' figures.  A panel is a small framed
% box holding one fixed layout of inputs, each drawn as a circle (+1) or a cross
% (-1); the mandatory argument is one letter per input, `o' or `x', in index
% order.  The layouts are the deck's own: \dichone one input, \dichtwo two,
% \dichthree three in general position, \dichcol three collinear ones and
% \dichfour four in convex position, so that the same coordinates recur in the
% big picture and in every panel beside it.  The optional argument `veto'
% strikes a panel through, marking a dichotomy that no line can produce.
% \dmark gives the same two marks at text size, for the dichotomy tables.
\tikzset{
	dichveto/.style={line width=2.2pt, draw=orange, line cap=round},
}
\newcommand{\dmarkat}[3]{%
	\ifthenelse{\equal{#3}{o}}%
	{\node[dpos] at (#1,#2) {};}%
	{\node[dneg] at (#1,#2) {\(\times\)};}}
\newcommand{\dmark}[1]{\tikz[baseline=-0.55ex]{\dmarkat{0}{0}{#1}}}
\newcommand{\dichonegrid}[1]{\dmarkat{1.00}{0.70}{#1}}
\newcommand{\dichtwogrid}[2]{\dmarkat{1.28}{0.98}{#1}\dmarkat{0.72}{0.46}{#2}}
\newcommand{\dichthreegrid}[3]{%
	\dmarkat{1.10}{1.08}{#1}\dmarkat{0.58}{0.70}{#2}\dmarkat{1.34}{0.32}{#3}}
\newcommand{\dichcolgrid}[3]{%
	\dmarkat{1.32}{1.08}{#1}\dmarkat{0.86}{0.70}{#2}\dmarkat{0.40}{0.32}{#3}}
\newcommand{\dichfourgrid}[4]{%
	\dmarkat{1.02}{1.10}{#1}\dmarkat{0.50}{0.70}{#2}%
	\dmarkat{1.28}{0.30}{#3}\dmarkat{1.62}{0.80}{#4}}
\newcommand{\dichframe}[3][]{%
	\begin{tikzpicture}[scale=0.62, baseline=(current bounding box.center)]
		\draw[plotbox] (0,0) rectangle (2.0,1.4);
		#2#3%
		\ifthenelse{\equal{#1}{veto}}{%
			\draw[dichveto] (0.12,0.12) -- (1.88,1.28);
			\draw[dichveto] (0.12,1.28) -- (1.88,0.12);}{}%
	\end{tikzpicture}}
\newcommand{\dichone}[2][]{\dichframe[#1]{\dichonegrid}{#2}}
\newcommand{\dichtwo}[2][]{\dichframe[#1]{\dichtwogrid}{#2}}
\newcommand{\dichthree}[2][]{\dichframe[#1]{\dichthreegrid}{#2}}
\newcommand{\dichcol}[2][]{\dichframe[#1]{\dichcolgrid}{#2}}
\newcommand{\dichfour}[2][]{\dichframe[#1]{\dichfourgrid}{#2}}

% The `bin' of marbles used all through the feasibility chapter.  `\binpath'
% traces the outline measured off the handout --- a waist just under the rim,
% a belly a third of the way up, a flat base.  It is a path macro rather than
% a style so that the very same coordinates can be \draw n (for the wall) and
% \clip ped (to keep stray marbles inside).  The bin is one unit tall and
% about 0.86 units wide, so `scale=3' in the enclosing scope gives a 3cm bin;
% marble nodes are placed in those scaled coordinates but keep an absolute
% `minimum size', which is why the diameter is passed explicitly.
\definecolor{marbleorange}{HTML}{F08000}
\definecolor{marblegreen}{HTML}{1B9E1B}
\newcommand{\binpath}{%
	plot[smooth] coordinates {(-0.340,1.000) (-0.299,0.795) (-0.326,0.631)
			(-0.395,0.467) (-0.432,0.303) (-0.416,0.139) (-0.363,0.000)}
	-- plot[smooth] coordinates {(0.363,0.000) (0.416,0.139) (0.432,0.303)
			(0.395,0.467) (0.326,0.631) (0.299,0.795) (0.340,1.000)}}
\tikzset{
	binwall/.style={line width=1pt, draw=black, line cap=round,
			line join=round},
	morange/.style={circle, fill=marbleorange, inner sep=0pt, minimum size=#1},
	morange/.default=2.4pt,
	mgreen/.style={circle, fill=marblegreen, inner sep=0pt, minimum size=#1},
	mgreen/.default=2.4pt,
	% the hand-drawn arrow that scoops a sample out of a bin
	scoop/.style={-{Stealth[length=2.6mm]}, line width=1pt, draw=black},
}

% Shading for the dichotomy tables of the bounding-function chapter.  A row
% whose first N-1 entries occur twice (once with each last entry) is a `pair',
% a row whose prefix occurs once is a `single'; the two shades of each
% alternate so that neighbouring pairs stay distinguishable.  `shatcell' marks
% the columns that a rejected dichotomy would shatter.
\colorlet{pairrow}{orange!30}
\colorlet{pairrowalt}{orange!14}
\colorlet{singlerow}{Purple!24}
\colorlet{singlerowalt}{Purple!11}
\colorlet{shatcell}{orange!35}

% One BAD region in the `space of data sets' pictures: a scribbled rosette
% centred at (x,y), in the given colour, turned by the optional angle so that
% neighbouring blobs do not look stamped.  Usage: \badblob[30]{0.9}{2.1}{red}.
\newcommand{\badblob}[4][0]{%
	\begin{scope}[shift={(#2,#3)}, rotate=#1]
		\draw[#4, line width=0.8pt, line join=round]
			plot[smooth, domain=0:1080, samples=110, variable=\t]
			({0.13*cos(\t) + 0.06*cos(3.7*\t)}, {0.13*sin(\t) - 0.06*sin(3.7*\t)});
	\end{scope}}

% The region of the B(N,k) and N^{k-1} tables (N >= 2, k >= 3) where the
% polynomial N^{k-1} already dominates B(N,k), in the VC-dimension chapter.
\colorlet{vcregion}{Purple!14}

% A tuning knob for the `degrees of freedom' picture: a dark disc inside a
% dial of 19 ticks labelled 0--18 counter-clockwise from the bottom, with a red
% pointer at the given tick.  Usage: \dialknob{x}{y}{tick}, inside a picture.
\tikzset{
	dialdisc/.style={circle, fill=MidnightBlue!55!gray, draw=black, thick,
			minimum size=10mm, inner sep=0pt},
	dialptr/.style={red, line width=1.1pt, line cap=round},
}
\newcommand{\dialknob}[3]{%
	\begin{scope}[shift={(#1,#2)}]
		\draw[black, thin] (0,0) circle (0.62);
		\foreach \i in {0,...,18}{%
			\pgfmathsetmacro{\a}{-90 + (\i + 0.5) * 360 / 19}%
			\draw[thin] (\a:0.62) -- (\a:0.72);
			\node[font=\fontsize{4}{4}\selectfont] at (\a:0.84) {\i};}
		\node[dialdisc] at (0,0) {};
		\pgfmathsetmacro{\a}{-90 + (#3 + 0.5) * 360 / 19}%
		\draw[dialptr] (0,0) -- (\a:0.55);
	\end{scope}}

% Cells of the 3x3 black/white pictures in the learning puzzle: `\puzzle{...}'
% takes nine digits, row by row, 1 for a black cell and 0 for a white one.
\newcommand{\puzzlecell}[3]{%
	\ifnum#3=1 \fill[black] (#1,#2) rectangle ++(1,1);\fi
	\draw[black, line width=0.4pt] (#1,#2) rectangle ++(1,1);}
\newcommand{\puzzle}[1]{%
	\begin{tikzpicture}[scale=0.32, baseline=(current bounding box.center)]
		\puzzlegrid#1\relax
	\end{tikzpicture}}
% ... and `\puzzleat{x}{y}{bits}' drops one grid with its lower-left corner at
% (x,y) inside a picture that is already open, one unit per cell, so that a
% whole board of them can be laid out with exact gaps.
\newcommand{\puzzleat}[3]{%
	\begin{scope}[shift={(#1,#2)}]\puzzlegrid#3\relax\end{scope}}
\newcommand{\puzzlegrid}[9]{%
	\puzzlecell{0}{2}{#1}\puzzlecell{1}{2}{#2}\puzzlecell{2}{2}{#3}%
	\puzzlecell{0}{1}{#4}\puzzlecell{1}{1}{#5}\puzzlecell{2}{1}{#6}%
	\puzzlecell{0}{0}{#7}\puzzlecell{1}{0}{#8}\puzzlecell{2}{0}{#9}}

% Pictures of the linear models (regression chapters onwards).
%   regpt / resid / regfit — a data point, its residual to the fitted
%       line/plane, and the fit itself, in the `illustration of linear
%       regression' plots;
%   gdcurve / gdpath — an error curve E_in(w) and the path of the iterates on
%       it, in the gradient-descent pictures;
%   spanfill — the column span of X in the hat-matrix pictures.
\tikzset{
	regpt/.style={circle, draw=black, thick, fill=white, inner sep=0pt,
			minimum size=4pt},
	resid/.style={red!85!black, semithick},
	regfit/.style={blue!85!black, thick},
	gdcurve/.style={Maroon, line width=1pt},
	gdpath/.style={black, thin, -{Stealth[length=1.6mm]}},
	spanfill/.style={fill=red!22},
}
% \thetacurve[scale]: the logistic function theta(s) = 1/(1+e^{-s}) on its own
% small axes, dashed asymptote at 1 --- it recurs beside every formula that
% uses it.
\newcommand{\thetacurve}[1][1]{%
	\begin{tikzpicture}[scale=#1, baseline=(current bounding box.center)]
		\draw[thin] (-1.9,0) -- (1.9,0);
		\draw[thin] (0,-0.25) -- (0,1.35);
		\draw[thin, dashed] (-1.9,1.1) -- (1.9,1.1);
		\node[left, font=\scriptsize] at (-1.9,1.1) {\(1\)};
		\node[left, font=\scriptsize] at (-1.9,0) {\(0\)};
		\node[below left, font=\scriptsize] at (1.9,0) {\(s\)};
		\draw[thick, blue] plot[domain=-1.9:1.9, samples=60]
			(\x, {1.1/(1 + exp(-2.4*\x))});
		\node[font=\scriptsize, text=blue] at (1.2,0.75) {\(\theta(s)\)};
	\end{tikzpicture}}
% \einbowl{label}: a smooth convex E_in(w) whose valley sits on the w axis,
% with the minimiser marked by a dot and the flat tangent there drawn as an
% arrow --- the picture of `gradient zero at the best w'.
\newcommand{\einbowl}[1]{%
	\begin{tikzpicture}[baseline=(current bounding box.center)]
		\draw[thin] (0,2.0) -- (0,0) -- (2.3,0);
		\node[left] at (0,1.05) {\(\Ein\)};
		\draw[gdcurve] plot[domain=0:2.25, samples=60]
			(\x, {1.9*((\x - 1.1)/1.1)^2});
		\draw[-{Stealth[length=1.6mm]}, thin] (1.1,0) -- (1.75,0);
		\fill (1.1,0) circle (1.5pt);
		\node[below] at (1.1,0) {#1};
	\end{tikzpicture}}
% \scoreunit{sign|id|theta}: the `linear scoring' diagram of a linear model ---
% inputs x_0, ..., x_d weighted into one score s, then the model's own
% transform of s (a step, the identity, or the logistic S) and out comes h(x).
\newcommand{\scoreunit}[1]{%
	\begin{tikzpicture}[font=\scriptsize, baseline=(current bounding box.center)]
		\foreach \lab/\yy in {x_0/0.75, x_1/0.45, x_2/0.05, x_d/-0.75}{%
			\node[left, inner sep=1pt] at (0,\yy) {\(\lab\)};
			\draw[thin] (0,\yy) -- (1.05,0);}
		\fill (1.0,0.15) -- (1.35,0) -- (1.0,-0.15) -- cycle;
		\draw (1.35,0) -- (1.8,0) node[midway, above] {\(s\)};
		\draw (2.1,0) circle (0.3);
		\ifthenelse{\equal{#1}{sign}}{%
			\draw[red, thick] (1.92,-0.13) -- (2.1,-0.13) -- (2.1,0.13)
				-- (2.28,0.13);}{}%
		\ifthenelse{\equal{#1}{id}}{%
			\draw[red, thick] (1.93,-0.17) -- (2.27,0.17);}{}%
		\ifthenelse{\equal{#1}{theta}}{%
			\draw[red, thick] plot[domain=-0.19:0.19, samples=30]
				({2.1 + \x}, {-0.15 + 0.3/(1 + exp(-30*\x))});}{}%
		\draw[-{Stealth[length=1.6mm]}] (2.4,0) -- (2.8,0)
			node[right, inner sep=1pt] {\(h(\bm{x})\)};
	\end{tikzpicture}}

% The running example of the nonlinear-transform chapter, measured off the
% handout: twenty inputs in [-1,1]^2, nine circles inside the circle
% x1^2 + x2^2 = 0.6 and eleven crosses outside it.  \circdataX draws them in
% X-space and \circdataZ their images z = (x1^2, x2^2) in Z-space, both in a
% picture whose unit is the data unit.  The first circle and the last cross are
% the examples the deck names x_1 and x_2; \circlabelsX and \circlabelsZ put
% those names (z_1, z_2 after the transform) beside them.  \xplotframe frames
% [-1,1]^2 with ticks at -1, 0, 1 and \zplotframe frames the deck's Z-space
% window [-0.05,1]^2 with ticks at 0, 0.5, 1.
\newcommand{\circdataX}{%
	\foreach \p in {(-0.494,0.363), (-0.311,-0.101), (-0.189,0.717),
			(-0.009,-0.173), (-0.007,0.374), (0.002,0.138), (0.085,0.322),
			(0.142,0.568), (0.171,-0.302)}{\node[dpos] at \p {};}%
	\foreach \p in {(-0.893,-0.946), (-0.803,0.877), (-0.774,0.551),
			(-0.748,-0.853), (-0.721,-0.720), (-0.647,0.772), (-0.634,-0.905),
			(0.491,0.920), (0.519,-0.715), (0.724,-0.794), (0.944,0.801)}%
		{\node[dneg] at \p {\(\times\)};}}
\newcommand{\circdataZ}{%
	\foreach \p in {(0.244,0.132), (0.097,0.010), (0.036,0.514),
			(0.000,0.030), (0.000,0.140), (0.000,0.019), (0.007,0.104),
			(0.020,0.323), (0.029,0.091)}{\node[dpos] at \p {};}%
	\foreach \p in {(0.797,0.895), (0.645,0.769), (0.599,0.304),
			(0.560,0.728), (0.520,0.518), (0.419,0.596), (0.402,0.819),
			(0.241,0.846), (0.269,0.511), (0.524,0.630), (0.891,0.642)}%
		{\node[dneg] at \p {\(\times\)};}}
\newcommand{\circlabelsX}{%
	\node[font=\scriptsize, below left, inner sep=1pt] at (-0.494,0.363)
		{\(\bm{x}_1\)};%
	\node[font=\scriptsize, below, inner sep=4pt] at (0.944,0.801)
		{\(\bm{x}_2\)};}
\newcommand{\circlabelsZ}{%
	\node[font=\scriptsize, below, inner sep=2pt] at (0.244,0.132)
		{\(\bm{z}_1\)};%
	\node[font=\scriptsize, below, inner sep=4pt] at (0.891,0.642)
		{\(\bm{z}_2\)};}
\newcommand{\xplotframe}{%
	\draw[plotbox] (-1,-1) rectangle (1,1);
	\foreach \t in {-1,0,1}{%
		\draw[plotbox] (\t,-1) -- (\t,-0.96);
		\draw[plotbox] (-1,\t) -- (-0.96,\t);
		\node[font=\tiny, below, inner sep=3pt] at (\t,-1) {\(\t\)};
		\node[font=\tiny, left, inner sep=1.5pt] at (-1,\t) {\(\t\)};}}
\newcommand{\zplotframe}{%
	\draw[plotbox] (-0.05,-0.05) rectangle (1,1);
	\foreach \t in {0,0.5,1}{%
		\draw[plotbox] (\t,-0.05) -- (\t,-0.03);
		\draw[plotbox] (-0.05,\t) -- (-0.03,\t);
		\node[font=\tiny, below, inner sep=3pt] at (\t,-0.05) {\(\t\)};
		\node[font=\tiny, left, inner sep=1.5pt] at (-0.05,\t) {\(\t\)};}}

% \vccurve: the `THE VC message' picture of the VC-dimension chapter --- in-sample
% error falling, model complexity rising and their sum, the out-of-sample
% bound, smallest at d_VC^* --- for chapters that point back at it.
\newcommand{\vccurve}{%
	\begin{tikzpicture}[baseline=(current bounding box.center)]
		\begin{axis}[
			width=7.0cm, height=4.8cm,
			xmin=0, xmax=10, ymin=0, ymax=10,
			axis lines=left, xtick={3.1}, xticklabels={\(d_{\mathrm{VC}}^*\)},
			ytick=\empty, xlabel={VC dimension, \(\dvc\)},
			ylabel={Error}, label style={font=\small},
			clip=false]
			\addplot[blue!80!black, thick, domain=0.3:10, samples=80]
				{6.5*exp(-0.55*x) + 0.4};
			\addplot[red!80!black, thick, domain=0:10, samples=40]
				{0.75*x};
			\addplot[Purple, thick, domain=0.3:10, samples=80]
				{6.5*exp(-0.55*x) + 0.4 + 0.75*x};
			\draw[dashed, Purple] (axis cs:3.1,0) -- (axis cs:3.1,4.54);
			\node[font=\scriptsize, text=blue!80!black, anchor=west]
				at (axis cs:6.2,1.1) {in-sample error};
			\node[font=\scriptsize, text=red!80!black, anchor=north west]
				at (axis cs:6.0,4.3) {model complexity};
			\node[font=\scriptsize, text=Purple, anchor=west]
				at (axis cs:6.5,9.7) {out-of-sample error};
		\end{axis}
	\end{tikzpicture}}

% \nestedH: the nested hypothesis sets H_0 c H_1 c H_2 c H_3 c ... of the
% structured-hypothesis-set pictures, as ellipses that share their leftmost
% point, each one swallowing the one before.
\newcommand{\nestedH}{%
	\begin{tikzpicture}[baseline=(current bounding box.center),
			font=\scriptsize]
		\foreach \a/\b in {0.50/0.22, 1.15/0.38, 1.85/0.52, 2.60/0.66}{%
			\draw[semithick] (\a,0) ellipse ({\a} and {\b});}
		\foreach \x/\lab in {0.50/0, 1.65/1, 3.00/2, 4.45/3}{%
			\node at (\x,0) {\(\mathcal{H}_{\lab}\)};}
		\node at (5.65,0) {\(\cdots\)};
	\end{tikzpicture}}

% Algorithm Env
\usepackage[linesnumbered,lined,vlined,ruled,commentsnumbered,resetcount,algochapter]{algorithm2e}
\SetKwComment{Comment}{// }{}
\SetArgSty{textsl}
\def\algocflineautorefname{Algorithm}
\counterwithin{algocfline}{chapter}

\makeatletter
\def\itemautorefname{\@gobble}
\makeatother

% ── Theorem infrastructure ────────────────────────────────────────────────
\usepackage{thmtools}

% mdframed — for Note-only environments
\usepackage[framemethod=TikZ]{mdframed}
\mdfsetup{skipabove=1em,skipbelow=0em}

\theoremstyle{definition}

% ── One look for every boxed environment ─────────────────────────────────
% A 3pt bar down the left in the environment's own colour, a 5% tint of that
% colour behind it, and no other rules.  mdframed splits across a page break
% correctly under XeTeX; tcolorbox does not (see the `splitsafe' note below),
% which is why none of these environments use tcolorbox any more.
\mdfdefinestyle{notebar}{
	linewidth=3pt,
	topline=false, bottomline=false, rightline=false,
	innerleftmargin=10pt, innerrightmargin=10pt,
	innertopmargin=6pt, innerbottommargin=8pt,
	skipabove=10pt, skipbelow=12pt,
	nobreak=false
}

% Results — head on its own line, one colour each so they are told apart.
\declaretheoremstyle[
	headfont=\sffamily\bfseries, bodyfont=\normalfont,
	headformat={\textcolor{RoyalBlue!70!black}{\NAME\ \NUMBER\NOTE}},
	headpunct={\\[3pt]}, postheadspace={0pt},
	mdframed={style=notebar, linecolor=RoyalBlue, backgroundcolor=RoyalBlue!5}
]{thmnapkinthm}

\declaretheoremstyle[
	headfont=\sffamily\bfseries, bodyfont=\normalfont,
	headformat={\textcolor{CadetBlue!70!black}{\NAME\ \NUMBER\NOTE}},
	headpunct={\\[3pt]}, postheadspace={0pt},
	mdframed={style=notebar, linecolor=CadetBlue, backgroundcolor=CadetBlue!5}
]{thmnapkinlem}

\declaretheoremstyle[
	headfont=\sffamily\bfseries, bodyfont=\normalfont,
	headformat={\textcolor{RoyalBlue!70!black}{\NAME\ \NUMBER\NOTE}},
	headpunct={\\[3pt]}, postheadspace={0pt},
	mdframed={style=notebar, linecolor=RoyalBlue, backgroundcolor=RoyalBlue!5}
]{thmnapkinprop}

\declaretheoremstyle[
	headfont=\sffamily\bfseries, bodyfont=\normalfont,
	headformat={\textcolor{RoyalBlue!70!black}{\NAME\ \NUMBER\NOTE}},
	headpunct={\\[3pt]}, postheadspace={0pt},
	mdframed={style=notebar, linecolor=RoyalBlue, backgroundcolor=RoyalBlue!5}
]{thmnapkincor}


% Exercises and questions — head inline, smaller body.
\declaretheoremstyle[
	headfont=\bfseries, bodyfont=\normalfont\small,
	mdframed={style=notebar, linecolor=black, backgroundcolor=RedViolet!5!gray!5}
]{thmnapkinques}

% Commentary — head inline.

\declaretheoremstyle[
	headfont=\bfseries\sffamily, bodyfont=\normalfont,
	headformat={\textcolor{SeaGreen!70!black}{\NAME\ \NUMBER\NOTE}},
	mdframed={style=notebar, linecolor=SeaGreen, backgroundcolor=SeaGreen!5}
]{thmgreen3box}

\declaretheoremstyle[
	headfont=\bfseries\sffamily, bodyfont=\normalfont,
	headformat={\textcolor{WildStrawberry!70!black}{\NAME\ \NUMBER\NOTE}},
	mdframed={style=notebar, linecolor=WildStrawberry, backgroundcolor=WildStrawberry!5}
]{thmpinkbox}

\declaretheoremstyle[
	headfont=\bfseries\sffamily, bodyfont=\normalfont,
	headformat={\textcolor{MidnightBlue!70!black}{\NAME\ \NUMBER\NOTE}},
	mdframed={style=notebar, linecolor=MidnightBlue, backgroundcolor=MidnightBlue!5}
]{thmblue2box}

\declaretheoremstyle[
	headfont=\bfseries\sffamily, bodyfont=\normalfont,
	headformat={\textcolor{Gray!70!black}{\NAME\ \NUMBER\NOTE}},
	mdframed={style=notebar, linecolor=Gray, backgroundcolor=Gray!5}
]{notgraybox}

\declaretheoremstyle[
	headfont=\bfseries\sffamily, bodyfont=\normalfont,
	headformat={\textcolor{Gray!70!black}{\NAME\ \NUMBER\NOTE}},
	mdframed={style=notebar, linecolor=Gray, backgroundcolor=Gray!5}
]{notgrayline}

\declaretheoremstyle[
	headfont=\bfseries\sffamily, bodyfont=\normalfont,
	headformat={\textcolor{Violet!70!black}{\NAME\ \NUMBER\NOTE}},
	mdframed={style=notebar, linecolor=Violet, backgroundcolor=Violet!5}
]{conjpurplebox}

% tcolorbox — for Napkin-style environments
\usepackage{tcolorbox}
\tcbuselibrary{breakable,skins,hooks}

% Patch tcolorbox so text can continue after a non-breakable box
\makeatletter
\tcbset{
	after app={%
		\ifx\tcb@drawcolorbox\tcb@drawcolorbox@breakable
		\else
			\@endparenv
		\fi
	}
}
\appto\tcb@use@after@lastbox{\@endparenv\@doendpe}
\makeatother

% `splitsafe' — use instead of `breakable' on any tcolorbox with a background.
%
% tcolorbox splits a breakable box across a page with \vsplit.  XeTeX carries
% the current colour in a stack driven by \special{color push/pop}, and that
% stack does not survive \vsplit: the continuation segment's body text is drawn
% in a leftover colour, in practice invisibly — the text is still in the PDF and
% pdftotext finds it, but nothing is readable.
%
% Re-checked on TeX Live 2024: still broken.  The trigger is a \color left open
% at the split point.  This is why the environments above are mdframed (which
% splits correctly) and colour their heads with \textcolor inside `headformat'
% rather than \color inside `headfont'.  `enhanced jigsaw', `opacityback=0' and
% \DocumentMetadata{} were all tried and none of them fix it.  LuaTeX's native
% \pdfcolorstack is split-safe, so the same source is fine there — switching
% engine is the only way to get these boxes to split.
%
% Keep them whole under XeTeX instead; they move to the next page as a unit.
\ifXeTeX
	\tcbset{splitsafe/.style={unbreakable}}
\else
	\tcbset{splitsafe/.style={breakable}}
\fi



% ── Boxed environments ───────────────────────────────────────────────────
\declaretheorem[style=thmnapkinthm, name=Theorem, numberwithin=section]{theorem}
\declaretheorem[style=thmnapkinlem, name=Lemma, sibling=theorem]{lemma}
\declaretheorem[style=thmnapkinprop, name=Proposition, sibling=theorem]{proposition}
\declaretheorem[style=thmnapkincor, name=Corollary, sibling=theorem]{corollary}

\declaretheorem[style=definition, name=Example, sibling=theorem]{eg}

\declaretheorem[style=thmnapkinques, name=Exercise, sibling=theorem]{exercise}
\declaretheorem[style=thmnapkinques, name=Question, sibling=theorem]{question}

\theoremstyle{definition}
\newtheorem{definition}[theorem]{Definition}
\newtheorem{fact}[theorem]{Fact}
\newtheorem{abuse}[theorem]{Abuse of Notation}

% ── Note-only environments (mdframed) ────────────────────────────────────
\declaretheorem[style=thmpinkbox, name=Problem, numberwithin=section]{problem}
\declaretheorem[style=thmpinkbox, name=Problem, numbered=no]{problem*}
\declaretheorem[style=thmblue2box, name=Claim, numberwithin=section]{claim}
\declaretheorem[style=conjpurplebox, name=Conjecture, numberwithin=section]{conjecture}

\declaretheorem[style=definition, sibling=theorem, name=Remark]{remark}
\declaretheorem[style=definition, sibling=theorem, name=Note]{note}
\declaretheorem[style=notgrayline, numbered=no, name=As previously seen]{prev}
\declaretheorem[style=notgrayline, numbered=no, name=Recall]{recall}
\declaretheorem[style=notgrayline, numbered=no, name=Observation]{observation}
\declaretheorem[style=thmgreen3box, numbered=no, name=Intuition]{intuition}
\declaretheorem[style=notgraybox, numbered=no, name=Notation]{notation}

\renewcommand\theHtheorem{\thesection.\arabic{theorem}}
\renewcommand\theHproblem{\thesection.\arabic{problem}}
\renewcommand\theHclaim{\thesection.\arabic{claim}}
\renewcommand\theHconjecture{\thesection.\arabic{conjecture}}

% ── Standard proof with black square ─────────────────────────────────────
\renewcommand{\qedsymbol}{$\blacksquare$}

% Napkin: subproof (nested proof, ■ distinguishes from outer)
\newenvironment{subproof}[1][Proof]{%
	\begin{proof}[#1]\renewcommand{\qedsymbol}{$\blacksquare$}}%
	{\end{proof}}

% Napkin: blue bold vocabulary term
\newcommand{\vocab}[1]{\textbf{\color{blue} #1}}

% Napkin: horizontal rule divider
\newcommand{\hrulebar}{%
	\par\hspace{\fill}\rule{0.95\linewidth}{.7pt}\hspace{\fill}%
	\par\nointerlineskip\vspace{\baselineskip}%
}

% Source marker: \source{slides p1}
\newcommand{\source}[1]{%
	\par\vspace{0.5em}%
	\noindent
	{\color{Gray}\rule[0.55ex]{1.5em}{0.6pt}}\enspace
	{\small\sffamily\bfseries\color{Gray!70!black}#1}\enspace
	{\color{Gray}\leaders\hrule height 0.6pt\hfill\kern0pt}%
	\par\nobreak\vspace{0.3em}%
}

% Napkin: green boxed moral / key takeaway
\newenvironment{moral}{%
	\begin{tcolorbox}[
		boxrule=0.4pt, colframe=green!70!black, sharp corners,
		standard jigsaw, opacityback=0,
		left=10pt, right=10pt, top=3pt, bottom=3pt,
		before skip=10pt, after skip=20pt]%
	\bfseries\color{green!50!black}}%
	{\end{tcolorbox}}

\usepackage{etoolbox}

% ── Page layout (KOMA-Script scrlayer-scrpage) ────────────────────────────
\usepackage[headsepline]{scrlayer-scrpage}
\renewcommand{\headfont}{}

\makeatletter

% Section/subsection: purple §, roman numeral subsections
\renewcommand*{\sectionformat}{\color{purple}{\lmrfont§}\thesection\autodot\enskip}
\renewcommand*{\subsectionformat}{\color{purple}{\lmrfont§}\thesubsection\autodot\enskip}
\renewcommand{\thesubsection}{\thesection.\roman{subsection}}

% Chapter: large purple number
\addtokomafont{chapterprefix}{\raggedleft}
\RedeclareSectionCommand[beforeskip=0.5em]{chapter}
\renewcommand*{\chapterformat}{%
	\mbox{\scalebox{1.5}{\chapappifchapterprefix{\nobreakspace}}%
	\scalebox{2.718}{\color{purple}\thechapter\autodot}\enskip}}

% Headers: chapter mark left, page number right
\automark[section]{chapter}
\ihead{\footnotesize\leftmark}
\ohead{\footnotesize\thepage}
\chead{}
\ifoot{}\ofoot{}\cfoot{}

\def\twodigits#1{\ifnum#1<10 0\fi\the#1}
\def\dateparts#1 #2 #3 #4 #5\relax{
	\marginpar{\footnotesize\textsf{\raggedright #1 #2 #3 #5}}
}

% KOMA-Script >= 3.31 makes a starred heading (\section*, \subsection*, ...)
% reset every counter numbered within it. Here unnumbered headings -- \lecture
% and in-text \section* -- sit in the middle of a numbered section, so the
% counters are saved when the starred branch starts and restored right after
% KOMA's reset, at the beginning of the heading's group. (KOMA calls each hook
% with one trailing argument, hence the \@gobble.)
\def\@starhead@restore{}
\def\@starhead@save#1{%
	\xdef\@starhead@restore{\unexpanded\expandafter{\@starhead@restore}\noexpand\setcounter{#1}{\the\value{#1}}}%
	\begingroup\let\@elt\@starhead@save\csname cl@#1\endcsname\endgroup
}
\def\@starhead@keepcounters#1{%
	\AddtoDoHook{heading/branch/star/#1}{%
		\gdef\@starhead@restore{}\@starhead@save{#1}\@gobble}%
	\AddtoDoHook{heading/begingroup/#1}{%
		\@starhead@restore\global\let\@starhead@restore\@empty\@gobble}%
}
\@starhead@keepcounters{section}
\@starhead@keepcounters{subsection}
\@starhead@keepcounters{subsubsection}
\def\@lecture{}%
\newcommand{\lecture}[3]{
	\ifthenelse{\isempty{#3}}{%
		\global\def\@lecture{Lecture #1}%
	}{%
		\global\def\@lecture{Lecture #1: #3}%
	}%
	\section*{\@lecture}
	\marginpar{\footnotesize\textsf{\raggedright #2}}
}

\usepackage{pgffor}
\newcommand{\lec}[2]{%
	\foreach \c in {#1,...,#2}{%
		\IfFileExists{Lectures/lec_\c.tex}{%
			\input{Lectures/lec_\c.tex}%
		}{}%
	}%
}

\makeatother

% notes
\usepackage[color=pink]{todonotes}
\usepackage{marginnote}
\let\marginpar\marginnote

% Appendix
\usepackage{appendix}
\def\chapterautorefname{Section}
\def\sectionautorefname{Section}
\def\appendixautorefname{Appendix}
\renewcommand\appendixname{Appendix}
\renewcommand\appendixtocname{Appendix}
\renewcommand\appendixpagename{Appendix}

\makeatletter
\patchcmd{\hyper@makecurrent}{%
	\ifx\Hy@param\Hy@chapterstring
		\let\Hy@param\Hy@chapapp
	\fi
}{%
	\iftoggle{inappendix}{%
		\@checkappendixparam{chapter}%
		\@checkappendixparam{section}%
		\@checkappendixparam{subsection}%
		\@checkappendixparam{subsubsection}%
		\@checkappendixparam{paragraph}%
		\@checkappendixparam{subparagraph}%
	}{}%
}{}{\errmessage{failed to patch}}

\newcommand*{\@checkappendixparam}[1]{%
	\def\@checkappendixparamtmp{#1}%
	\ifx\Hy@param\@checkappendixparamtmp
		\let\Hy@param\Hy@appendixstring
	\fi
}
\makeatother

\newtoggle{inappendix}
\togglefalse{inappendix}

\apptocmd{\appendix}{\toggletrue{inappendix}}{}{\errmessage{failed to patch}}
\apptocmd{\subappendices}{\toggletrue{inappendix}}{}{\errmessage{failed to patch}}

\setcounter{tocdepth}{1}

% Room for three-digit page numbers in the table of contents (the bold chapter
% entries overflow the default width from page 100 on).
\makeatletter
\renewcommand*{\@pnumwidth}{2.2em}
\makeatother
